% 第1章 1.4.2：Young 不等式的面积解释，p=3，q=3/2。
% 曲线 y=t^{p-1}=t^2；a=1，b=0.6。蓝色区域面积 a^p/p=1/3，绿色区域面积 b^q/q=(2/3)(0.6)^{3/2}≈0.310；
% 二者覆盖矩形 [0,a]x[0,b]（面积 ab=0.6），0.6<=0.643。
\begin{tikzpicture}[font=\figurefont\footnotesize]
\begin{axis}[width=56mm,height=48mm,axis lines=left,
  x axis line style={draw=FigureStroke,line width=.6pt,-{Stealth[length=2mm]}},
  y axis line style={draw=FigureStroke,line width=.6pt,-{Stealth[length=2mm]}},
  tick style={draw=FigureStroke},tick label style={font=\figurefont\scriptsize,text=FigureInk},
  xmin=0,xmax=1.25,ymin=0,ymax=1.2,xtick={1},xticklabels={$a$},ytick={0.6},yticklabels={$b$},
  xlabel={$t$},ylabel={$s$},xlabel style={at={(axis description cs:1,0)},anchor=north west,yshift=1mm,font=\figurefont\scriptsize},
  ylabel style={at={(axis description cs:0,1)},anchor=south east,rotate=-90,font=\figurefont\scriptsize},samples=100]
\addplot[fill=FigureBlueFill,draw=none,domain=0:1] {x^2} \closedcycle;
\addplot[fill=FigureGreenFill,draw=none,domain=0:0.6,variable=\s] ({sqrt(\s)},{\s}) -- (axis cs:0,0.6) -- (axis cs:0,0) -- cycle;
\addplot[draw=FigureStroke,line width=.75pt] coordinates {(0,0.6) (1,0.6) (1,0)};
\addplot[draw=FigureBlue,line width=1.1pt,domain=0:1.08] {x^2};
\node[font=\figurefont\scriptsize,text=FigureBlue,anchor=east] at (axis cs:0.97,1.02) {$s=t^{p-1}$};
\node[font=\figurefont\scriptsize,text=FigureBlue] at (axis cs:0.80,0.22) {$\dfrac{a^p}{p}$};
\node[font=\figurefont\scriptsize,text=FigureGreen] at (axis cs:0.24,0.40) {$\dfrac{b^q}{q}$};
\node[font=\figurefont\scriptsize,text=FigureInk,anchor=south] at (axis cs:0.5,0.6) {矩形面积 $ab$};
\end{axis}
\end{tikzpicture}
